\section{Struktur dan Operasi Symbol Table}
\label{sec:ch15-struktur}

\subsection{Atribut Simbol}

Tabel simbol memetakan nama ke informasi tentang nama tersebut.
Atribut yang umum:
\begin{itemize}[leftmargin=*]
  \item nama dan panjangnya;
  \item jenis simbol: variabel, fungsi, atau parameter;
  \item tipe;
  \item scope asal;
  \item lokasi deklarasi (baris, kolom);
  \item alamat memori atau indeks register (untuk fase berikutnya).
\end{itemize}

\begin{definisi}[Tabel simbol]
Tabel simbol adalah struktur data yang mendukung operasi
\code{define}, \code{lookup}, dan \code{scope management} untuk
seluruh identifier dalam program.
\end{definisi}

\subsection{Operasi Dasar}

\begin{table}[H]
\centering
\caption{Operasi tabel simbol}
\label{tab:ch15-operasi}
\begin{tabularx}{\textwidth}{@{}l X@{}}
\toprule
\textbf{Operasi} & \textbf{Deskripsi} \\
\midrule
\code{define} & Menambahkan simbol ke scope saat ini \\
\code{lookup} & Mencari simbol dari scope dalam ke luar \\
\code{lookup\_current} & Mencari simbol hanya di scope saat ini \\
\code{enter\_scope} & Mendorong scope baru \\
\code{leave\_scope} & Melepas scope teratas \\
\bottomrule
\end{tabularx}
\end{table}

\subsection{Representasi}

Tabel simbol dapat diimplementasikan sebagai:
\begin{itemize}[leftmargin=*]
  \item daftar tertaut sederhana;
  \item hash table dengan chaining;
  \item pohon pencarian seimbang;
  \item kombinasi hash table dan stack scope.
\end{itemize}

Untuk proyek Subset-C, setiap scope memiliki daftar simbol
tertaut. Implementasi ini sederhana dan cukup untuk program
perkuliahan. Untuk program besar, hash table diperlukan agar
\code{lookup} tetap cepat.

\begin{figure}[H]
\centering
\begin{tikzpicture}[node distance=0mm]
  \node[memorycell, minimum width=3.2cm] (s0) {Scope global};
  \node[memorycell, minimum width=3.2cm, below=of s0] (s1)
        {Scope fungsi};
  \node[memorycell, minimum width=3.2cm, below=of s1] (s2)
        {Scope blok};
  \draw[arrow] (s2) -- (s1);
  \draw[arrow] (s1) -- (s0);
\end{tikzpicture}
\caption{Stack scope: pencarian bergerak dari dalam ke luar}
\label{fig:ch15-scope-stack}
\end{figure}

\subsection{Tabel Informasi}

Guntoro membedakan beberapa tabel informasi yang sering dipakai
compiler. Tabel-tabel ini memperjelas atribut apa saja yang perlu
disimpan di samping nama simbol.

\begin{table}[H]
\centering
\caption{Jenis tabel informasi dan isinya}
\label{tab:ch15-informasi}
\begin{tabularx}{\textwidth}{@{}l X@{}}
\toprule
\textbf{Tabel} & \textbf{Informasi yang disimpan} \\
\midrule
Identifier & Semua identifier yang muncul, beserta nomor urut dan
atributnya \\
Array & Dimensi, tipe elemen, batas bawah dan atas, serta ukuran
tiap dimensi \\
Blok & Variabel yang dideklarasikan dalam blok yang sama \\
Real & Elemen bernilai real, termasuk presisi dan representasinya \\
String & Panjang, isi, dan lokasi penyimpanan string \\
Display & Blok yang sedang aktif pada scope bersarang \\
\bottomrule
\end{tabularx}
\end{table}

Tabel display berguna saat lookup harus mengetahui blok aktif tanpa
menelusuri seluruh stack scope. Pada proyek Subset-C, tabel simbol
per scope sudah cukup; tabel informasi di atas menjadi acuan bila
bahasa diperluas dengan array, string, atau blok bertingkat.

\subsection{Resolusi Nama}

Saat \code{lookup} dipanggil, tabel mencari di scope saat ini. Jika
tidak ditemukan, pencarian berlanjut ke scope induk, lalu ke scope
global. Jika tetap tidak ditemukan, identifier dianggap tidak
dideklarasikan.

Urutan ini penting untuk shadowing: identifier di scope dalam
menutupi identifier dengan nama sama di scope luar. Shadowing
diizinkan, sedangkan deklarasi ganda dalam scope yang sama merupakan
kesalahan.
